\chapter{イントロダクション}

現代の大学の数学科の数学教育では
数学の歴史的な発展を学ぶことは重視されておらず，
しかも，遠近法や射影幾何学や一次分数変換や非ユークリッド幾何学などについては説明すらなされないまま，
一般位相や関数論，多様体論，位相幾何学，代数を学び，
その応用の中で突然射影空間，射影変換，楕円曲線，保形形式が現れるような構成になっている．
時間の関係でこのような形になるのは仕方が無いのかもしれないが，
そのために数学に挫折してしまう学生の多いことは悲しいことである．
そこでこの本ではユークリッド幾何学の完成以降の
現代数学に繋がるユークリッド幾何学の発展について解説し，
数学科で学ぶ数学へのイントロダクションになることを期待している．
具体的には，$\SO(3)$や$\Ot(3)$の有限群，
ユークリッド平面の装飾群の(特にココンパクトなもの17個の)分類と
完備平坦リーマン多様体の分類を与える．
これは群作用の観点からまとめてあるので群論の入門になると思うが，
リーマン幾何学の入門ともなっている．
そのあと遠近法や3DCG，射影幾何学，
実射影空間，リーマン球面$\mathbf CP^1$と一次変換の関係(分類)，
トーラスのモジュライと上半平面の双曲幾何，楕円曲線とモジュラー関数
などを紹介しており，関数論との関係も強調しておきたい．
必ずしも歴史的に正確なことにはこだわらないが，
おおよその流れはおえるような形で記述を与えたい．

数学的な予備知識については数学科の大学2年生ぐらいまでで習うことは仮定しているが，
逆にこの本の内容を深く理解しようという視点からそれらの教科書を復習して
もらえれば幸いである．

大学の数学を始めると絵や図を描いて説明してはいけないなどの不思議なルールが
課せられることがあるが，それは初学者への方便であり，
絵や図を使った方が分かり易いところでは沢山の絵や図を入れる方針で本文は書いた．
ホームページにこの本の図のソースファイルや関連した動画やアプリを
\href{https://snow00two.github.io/mySite/}{私のサイト}に用意したので
それらを動かしながら理解を深めてもらいたい．